\subsection{希尔伯特空间的正交子空间的希尔伯特和}

定义 2. --- 称希尔伯特空间 E 是其一族闭向量子空间 \((\mathbf{E}_{i})_{i\in I}\) 的希尔伯特和，如果：

\begin{enumerate}
\def\labelenumi{\arabic{enumi})}
\item
  对 I 中两个不同的指标 \(i\)、\(k\)，子空间 \(\mathrm{E}_i\) 与 \(\mathrm{E}_k\) 在 E 中正交；
\item
  由诸 \(\mathbf{E}_i\) 的并生成的闭向量子空间就是 E。
\end{enumerate}

定理 1. --- 设 E 为希尔伯特空间，它是一族闭向量子空间 \((\mathrm{E}_{i})_{i\in\mathbf{I}}\) 的希尔伯特和。存在唯一的、从 E 到 \(\bigoplus_{i\in I}E_{i}=F\)（族 \((\mathrm{E}_{i})\) 的外希尔伯特和）上的同构 f，使得对所有 \(i\in I\)，\(f\) 在 \(\mathbf{E}\) 上的限制就是典范映射 \(f_{i}\)（从 \(\mathbf{E}_i\) 到 \(\mathbf{F}\) 中）。

设 \(S \subset F\) 为诸 \(E_i\) 的 « 代数 » 直和，并设 \(g\) 为线性映射 \((x_i)_{i \in I} \mapsto \sum_{i \in I} x_i\)（从 \(S\) 到 \(E\) 中）。我们将证明 \(g\) 是从准希尔伯特空间 S 到（准希尔伯特）子空间 \(g(S)\)（E 的由诸 \(E_{i}\) 的并生成的子空间）上的同构：因为，对两个元素 \(\mathbf{x} = (x_{i})_{i \in \mathrm{I}}\)、\(\mathbf{y} = (y_{i})_{i \in \mathrm{I}}\)，我们有

\[
\langle g (\mathbf {x}) | g (\mathbf {y}) \rangle = \left\langle \sum_ {i \in \mathrm{I}} x _ {i} \right| \sum_ {i \in \mathrm{I}} y _ {i} \rangle = \sum_ {(i, k) \in \mathrm{I} \times \mathrm{I}} \left\langle x _ {i} \mid y _ {k} \right\rangle .
\]

但若 \(i \neq k\)，则由假设 \(\langle x_i | y_k \rangle = 0\)，因而

\[
\langle g (\mathbf {x}) | g (\mathbf {y}) \rangle = \sum_ {i \in \mathbf {I}} \langle x _ {i} | y _ {i} \rangle = \langle \mathbf {x} | \mathbf {y} \rangle ;
\]

这就证明了我们的断言。由于 S 在 F 中稠密且 \(g(S)\) 在 E 中稠密，同构 g 延拓为从 F 到 E 上的同构 \(\overline{g}\)（V, p.~8, 推论）。显然 \(\overline{g}\) 的逆同构 f 就是所求的映射；其唯一性由以下事实得出：E 中由诸 \(E_{i}\) 的并生成的闭子空间就是 E 自身。

当 E 是一族子空间 \((\mathrm{E}_{i})_{i\in I}\) 的希尔伯特和时，我们常借助同构 f 把 E 与诸 \(E_{i}\) 的外希尔伯特和 F 等同起来。~若集 I 为有限，则说 E 是族 \((\mathrm{E}_{i})_{i\in I}\) 的希尔伯特和，就意味着诸 \(E_{i}\) 两两正交，且向量空间 E 是子空间族 \((\mathrm{E}_{i})_{i\in I}\) 的直和。

推论 1. --- 设 E 为希尔伯特空间，它是一族闭向量子空间 \((\mathbf{E}_{i})_{i\in\mathbf{I}}\) 的希尔伯特和；对所有 \(i\in I\)，设 \(p_{E_{i}}\) 为从 E 到 \(E_{i}\) 上的正交投影算子（V, p.~13）。

\begin{enumerate}
\def\labelenumi{\alph{enumi})}
\tightlist
\item
  对所有 \(x \in \mathbb{E}\)，族 \((\| p_{\mathbf{E}_i}(x)\|^{2})_{i\in \mathbf{I}}\) 在 \(\mathbf{R}\) 中可和，族 \((p_{\mathbf{E}_i}(x))_{i\in \mathbf{I}}\) 在 \(\mathbb{E}\) 中可和，并且我们有
\end{enumerate}

\[
\left\| x \right\| ^ {2} = \sum_ {i \in \mathbf {I}} \left\| p _ {\mathrm{E} _ {i}} (x) \right\| ^ {2}, \quad x = \sum_ {i \in \mathbf {I}} p _ {\mathrm{E} _ {i}} (x).
\]

\begin{enumerate}
\def\labelenumi{\alph{enumi})}
\setcounter{enumi}{1}
\item
  反之，若 \((x_{i})_{i\in \mathbf{I}}\) 是 \(\mathbf{E}\) 中元素的一族，使得 \(x_{i}\in E_{i}\) 对所有 \(i\in \mathbf{I}\) 成立，且 \(\sum_{i\in \mathbf{I}}\| x_i\|^2 < + \infty\)，则此族可和，且其和 \(x\) 是 \(\mathbf{E}\) 中唯一的满足 \(p_{\mathrm{E}_i}(x) = x_i\)（对所有 \(i\in \mathbf{I}\)）的点。
\item
  对每一对点 \(x, y\)（属于 \(\mathbf{E}\)），我们有
\end{enumerate}

\[
\langle x | y \rangle = \sum_ {i \in \mathbf {I}} \left\langle p _ {\mathrm{E} _ {i}} (x) \mid p _ {\mathrm{E} _ {i}} (y) \right\rangle .
\]

这些性质对诸 \(\mathbf{E}_i\) 的外希尔伯特和而言事实上是显然的，并可经由同构转移到 \(\mathbf{E}\) 上。

推论 2. --- 设 E 为分离的准希尔伯特空间，\((\mathrm{E}_{i})_{i\in\mathrm{I}}\) 为 E 的一族完备向量子空间，使得对 I 中每一对不同的指标 i、k，子空间 \(E_{i}\) 与 \(E_{k}\) 正交。设 V 为 E 中由诸 \(E_{i}\) 的并生成的闭向量子空间。对每个 \(i\in I\)，设 \(p_{E_{i}}\) 为从 E 到 \(E_{i}\) 上的正交投影算子。设 \(x\in E\)。

\begin{enumerate}
\def\labelenumi{\arabic{enumi})}
\item
  我们有 \(\sum\left\|p_{E_{i}}(x)\right\|^{2}\leqslant\left\|x\right\|^{2}\)。
\item
  下列条件等价：a) \(x \in \mathbf{V}\)；b) \(\sum_{i \in I} \| p_{\mathrm{E}_i}(x) \|^2 = \| x \|^2\)；c) 族 \((p_{\mathrm{E}_i}(x))_{i \in I}\) 在 E 中可和，且 \(x = \sum_{i \in I} p_{\mathrm{E}_i}(x)\)。
\item
  假设 V 完备。则族 \((p_{\mathrm{E}_i}(x))_{i\in \mathbf{I}}\) 在 E 中可和，且
\end{enumerate}

\[
p _ {\mathbf {V}} (x) = \sum_ {i \in \mathbf {I}} p _ {\mathrm{E} _ {i}} (x), \quad \left\| p _ {\mathbf {V}} (x) \right\| ^ {2} = \sum_ {i \in \mathbf {I}} \left\| p _ {\mathrm{E} _ {i}} (x) \right\| ^ {2},
\]

其中 \(p_{V}\) 表示从 E 到 V 上的正交投影算子。

设 \(\hat{\mathbf{E}}\) 为 \(\mathbf{E}\) 的希尔伯特空间完备化；我们把 \(\mathbf{E}\) 与 \(\hat{\mathbf{E}}\) 的一个稠密子空间等同起来；诸 \(\mathbf{E}_i\) 因完备而是 \(\hat{\mathbf{E}}\) 的闭子空间。闭包 \(\overline{\mathbf{V}}\)（\(\mathbf{V}\) 在 \(\hat{\mathbf{E}}\) 中的闭包）是 \(\hat{\mathbf{E}}\) 中由诸 \(\mathbf{E}_i\) 的并生成的闭向量子空间，且 \(\mathbf{V} = \overline{\mathbf{V}} \cap \mathbf{E}\)。空间 \(\hat{\mathbf{E}}\) 是诸 \(\mathbf{E}_i\) 与子空间 \(\mathbf{W}\)（即 \(\overline{\mathbf{V}}\) 在 \(\hat{\mathbf{E}}\) 中的正交补）的希尔伯特和；令 \(x_0 = p_{\mathrm{W}}(x)\) 与 \(x_i = p_{\mathrm{E}_i}(x)\)（对所有 \(i \in \mathbf{I}\)）。由推论 1，我们有 \(\| x\|^2 = \| x_0\|^2 + \sum_{i \in \mathbf{I}} \| x_i\|^2\)，且 \(x = x_0 + \sum_{i \in \mathbf{I}} x_i\)（在 \(\hat{\mathbf{E}}\) 中）。由此即得断言 1)，以及 2) 的条件 \(b)\) 与 \(c)\) 等价于条件 \(x_0 = 0\)、从而等价于条件 \(x \in \mathbf{V}\) 这一事实。最后，若 \(\mathbf{V}\) 完备，并令 \(x' = p_{\mathrm{V}}(x)\)，我们有 \(x' - x_i = (x - x_i) - (x - p_{\mathrm{V}}(x))\)，因而 \(x' - x_i\) 与 \(\mathbf{E}_i\) 正交，于是 \(x_i = p_{\mathrm{E}_i}(x')\) 对所有 \(i \in \mathbf{I}\) 成立；此时只需对向量 \(x'\) 应用性质 2) 即可。

注. --- 设 E 为分离的准希尔伯特空间，\((\mathrm{V}_{i})_{i\in\mathrm{I}}\) 为 E 的一族向量子空间，使得对每一对不同的指标 i、k，子空间 \(V_{i}\) 与 \(V_{k}\) 正交。于是，对每个 \(k\in I\)，\(V_{k}\) 与闭向量子空间 \(W_{k}\)（由诸 \(V_{i}\)（\(i\neq k\)）的并生成）的交化为 0；因为，若 x 既属于 \(V_{k}\) 又属于 \(W_{k}\)，则它与所有 \(V_{i}\)（\(i\neq k\)）正交，从而与 \(W_{k}\) 正交。特别地，x 与自身正交，故为零。

命题 2. --- 设 E 为希尔伯特空间，\((\mathrm{V}_{\lambda})_{\lambda\in\mathrm{L}}\) 为 E 的一族闭向量子空间；对每个 \(\lambda\in L\)，设 \((\mathrm{W}_{\lambda\mu})_{\mu\in\mathrm{M}_{\lambda}}\) 为 \(V_{\lambda}\) 的一族闭向量子空间，使得 \(V_{\lambda}\) 是由此族之并生成的闭向量子空间。为使 E 是族 \((\mathrm{W}_{\lambda\mu})_{\lambda\in\mathrm{L},\mu\in\mathrm{M}_{\lambda}}\) 的希尔伯特和，必要且充分的是：E 是族 \((\mathrm{V}_{\lambda})_{\lambda\in\mathrm{L}}\) 的希尔伯特和，且对每个 \(\lambda\in L\)，\(V_{\lambda}\) 是族 \((\mathrm{W}_{\lambda\mu})_{\mu\in\mathrm{M}_{\lambda}}\) 的希尔伯特和（« 希尔伯特和的结合性 »）。为证条件必要，只需看出 \(V_{\alpha}\) 与 \(V_{\beta}\) 正交（若 \(\alpha\neq\beta\)）。但是，\(W_{\alpha\mu}\)（\(\mu\in M_{\alpha}\)）的每个元素与所有 \(W_{\beta\nu}\)（\(\nu\in M_{\beta}\)）正交，从而与它们所生成的闭向量子空间 \(V_{\beta}\) 正交；同样的论证进而表明 \(V_{\beta}\) 的每个元素与 \(V_{\alpha}\) 正交，因为它与所有 \(W_{\alpha\mu}\)（\(\mu\in M_{\alpha}\)）正交。

为证条件充分，只需验证：若条件满足，则 \(\mathrm{E}\) 等于闭向量子空间 \(\mathrm{F}\)（由诸 \(\mathbf{W}_{\lambda \mu}\)（\(\lambda \in \mathbf{L}\)，\(\mu\in M_{\lambda}\)）的并生成）；而对每个 \(\lambda\in L\)，F 包含由诸 \(W_{\lambda\mu}\)（其中 \(\mu\in M_{\lambda}\)）的并生成的闭向量子空间，也就是说 F 包含 \(V_{\lambda}\)；于是 F 是由诸 \(V_{\lambda}\) 的并生成的闭向量子空间，而由假设它就是 E。

